\title[The Optimal Contract Makes the Flexible Load Decarbonize (Week 3)]
{The Optimal Contract Makes a Flexible Computing Load Decarbonize the Grid:\\
A Shape-Independent Guarantee}



\begin{abstract}
This is the Role~1 Week~3 deliverable, and it is a positive resolution of the Week-2
carbon question. Week~2 showed that under the \emph{naive} dispatch rule---run the
flexible computing load (\textsc{fcl}) on spare fossil whenever it pays---a finite
interior optimum carries an irreducible fossil floor, so the carbon sign is ambiguous
and, in general, not even decidable from the equilibrium cutoffs. Here we change the
contract and the answer changes with it. We study the regulator-imposed
\emph{surplus contract}, which runs the \textsc{fcl} on renewable surplus and
curtails it in fossil hours, and which minimizes the load's fossil footprint
(driving it to zero). Three results follow, and all are positive. \emph{First}, under
the surplus contract the load burns \emph{zero} fossil, so the Week-2 floor is gone by
construction. \emph{Second}---the keystone---whenever the load is adopted the contract
strictly raises renewable investment, $\KRc>\KRo$, and total carbon \emph{strictly
falls}, $\E[\fossil^{c}]<\E[\fossil^{\circ}]$. The force of this is that, unlike the
naive rule, it needs \emph{no condition on the shape of $F$}: under the standing
regularity (the renewable problem is single-peaked) and adoption, carbon falls whatever
the tails or skew of demand. The carbon sign, which Week~2 showed is decided by the
distributional shape under the naive rule, is freed from that shape by the contract.
\emph{Third}, the contract is a clean Pareto improvement on the no-load benchmark:
profit strictly rises and carbon strictly falls, while reliability is held exactly at
the no-load level (the surplus load never touches fossil, so it leaves the blackout
target unchanged). The two standing hypotheses are adoption ($\CM<R\,F(0;\KRo)$, the
load is cheap enough to run on the available surplus) and the Week-2 single-peakedness;
neither is a \emph{sign-determining shape condition}---single-peakedness is the
regularity that guarantees a well-defined optimum, not a restriction that decides which
way carbon moves---and we are explicit that the surplus optimum may sit at the
renewable-capacity bound when the surplus channel is strong. \emph{Fourth}, we place the naive and surplus contracts as the two ends of a single
one-parameter family---grant the load a per-state fossil budget $\tau$, with $\tau=0$ the
surplus rule and $\tau\to\infty$ the naive rule. Profit rises with $\tau$ (proved), and
carbon-optimality of surplus-only reduces, by an \emph{exact identity}
$C(\tau)-C(0)=\Phi(\tau)(1-\Lambda(\tau))$, to a single \emph{leverage} condition
$\Lambda\le1$ (the reduction proved; the condition holds with room to spare, $\Lambda<0$,
in every case we scan): the contracts trace a monotone carbon--profit frontier with \emph{no free
lunch}. So Weeks~2 and~3 are not rival models but the carbon-worst and carbon-best corners
of one optimal-contract family, the regulator's reliability and participation requirements
selecting the point between. Throughout, the core decarbonization is
\emph{distributionally robust}: it holds for the worst-case demand law over the regularity
class, so the regulator needs no knowledge of $F$. Every identity is
checked against quadrature in \texttt{verify\_week3.py}. The message is clean: a flexible
computing load is a decarbonization instrument, and contract design is the lever that
makes it one.
\end{abstract}



%---------------------------------------------------------------------
\section{Setup: the model and the contract space}\label{r1w3-sec:setup}

We keep the model and notation of the Week-1 and Week-2 deliverables. A regulated
utility chooses renewable capacity $\KR$, dispatchable fossil $\KF$, and
flexible-computing-load (\textsc{fcl}) capacity $\KM$. Residual demand is
$\tilde x=\tilde D-\KR\tilde S\in[-\KR,\Dbar]$ with CDF $F(x;\KR)$, density
$f(x;\KR)>0$ on $(-\KR,\Dbar)$, and $\FKR\eqdef\partial F/\partial\KR>0$ on the
interior of the support (more renewables strictly shift residual demand toward surplus
in every interior state---true for the scale model $\tilde x=\tilde D-\KR\tilde S$,
where $\FKR=\E[\tilde S\,f_{\tilde D}(x+\KR\tilde S)]>0$ wherever $f_{\tilde D}>0$). This
positivity is the only structural property of $F$ we use, and it holds across surplus
states ($x<0$) and fossil states ($x>0$) alike. Per-MWh: \textsc{fcl} revenue $R$,
fossil cost $c$, price $p$, blackout penalty $\beta$; amortized capital costs
$\CR,\CF,\CM$.

\begin{assumption}[Regularity, as in Week 1]\label{r1w3-ass:reg}
$R>c+\CM$,\quad $R<c+\CF+\CM$,\quad $\bp>c+\CF$,\quad $\bp>R$.
\end{assumption}

\paragraph{Dispatch and fossil accounting.} In state $x$, consumers are served first
(they outrank the \textsc{fcl} for scarce fossil), so consumer fossil is
$\min\{(x)_+,\KF\}$. Renewable surplus $(-x)_+$ is free power. A \emph{dispatch
contract} is a measurable schedule $d:\mathbb R\to[0,\KM]$ with
$d(x)\le(\KF-x)_+$ (the \textsc{fcl} can draw surplus plus fossil headroom); the
fossil it burns is $(d(x)-(-x)_+)_+$ (surplus first, then fossil). Total fossil is
\begin{equation}\label{r1w3-eq:totfossil}
  \E[\fossil]=\underbrace{\E\!\big[\min\{(\tilde x)_+,\KF\}\big]}_{\text{consumer fossil}}
  \;+\;\underbrace{\E\!\big[(d(\tilde x)-(-\tilde x)_+)_+\big]}_{\text{\textsc{fcl} fossil}} .
\end{equation}
Consumer fossil equals $\int_0^{\KF}(1-F)\,dx$ and does not depend on the contract
$d$; only the second term does.

\paragraph{Two contracts.} The Week-2 analysis used the \emph{firm-optimal} (naive)
contract $d^{*}(x)=\min\{(\KF-x)_+,\KM\}$, which runs the load whenever the marginal
$R-c>0$ is positive---hence on fossil headroom too. We contrast it with the
\emph{surplus contract}
\begin{equation}\label{r1w3-eq:surpcontract}
  d^{c}(x)\eqdef\min\{(-x)_+,\KM\} ,
\end{equation}
which runs the load only on renewable surplus and curtails it whenever fossil would be
needed. We separate the two decision-makers cleanly: the \emph{regulator} sets the
real-time dispatch contract $d(\cdot)$ (the regulated utility dispatches as directed),
and the \emph{firm} then chooses capacities $(\KR,\KF,\KM)$ to maximize profit under
that contract. The surplus contract is regulator-imposed and is \emph{not} dispatch
incentive-compatible: with $R>c$ the firm would prefer to run the load on fossil
headroom too (the naive $d^{*}$). It is enforced as a curtailment rule, exactly as
demand-response and interruptibility tariffs are. The no-\textsc{fcl} benchmark
$(\KRo,\KFo)$ is the program with $\KM=0$, with fossil
$\E[\fossil^{\circ}]=\int_0^{\KFo}(1-F(x;\KRo))\,dx$ and the distribution-free
cutoff
\begin{equation}\label{r1w3-eq:phio}
  F(\KFo;\KRo)=\varphi^{\circ}\eqdef\frac{\bp-c-\CF}{\bp-c}.
\end{equation}

%=====================================================================
\section{The fossil-minimizing dispatch is surplus-only}\label{r1w3-sec:optcontract}
%=====================================================================

We first show that, at any fixed capacities, the surplus contract minimizes the load's
fossil footprint---driving it to zero---and that among zero-footprint contracts it is
the one that monetizes the most surplus, hence the natural carbon-oriented contract.
(We do \emph{not} claim global carbon-optimality over the bilevel game: a different
contract could induce a larger $\KR$ at the price of some \textsc{fcl} fossil. Week~2
showed that trade---the naive rule buys a bigger capacity bump but pays the fossil
floor---loses on net; the surplus contract takes the other side, zero footprint with a
strict capacity gain, and \S\ref{r1w3-sec:decarb} shows that side wins distribution-wide.)

\begin{proposition}[Fossil-minimizing dispatch]\label{r1w3-prop:surplusopt}
Fix capacities $(\KR,\KF,\KM)$. Among all admissible dispatch schedules $d$, total
fossil \eqref{r1w3-eq:totfossil} is minimized, state by state, by running the load on
surplus alone, $d^{c}(x)=\min\{(-x)_+,\KM\}$, at which the \textsc{fcl} fossil term is
identically zero. Any schedule that ever dispatches beyond the surplus,
$d(x)>(-x)_+$, burns strictly more fossil in that state.
\end{proposition}

\begin{proof}
Consumer fossil in \eqref{r1w3-eq:totfossil} is independent of $d$, so minimizing total
fossil is minimizing $\E[(d(\tilde x)-(-\tilde x)_+)_+]$ pointwise. The integrand
$(d(x)-(-x)_+)_+$ is zero whenever $d(x)\le(-x)_+$ and strictly increasing in $d(x)$
beyond that. Setting $d(x)=\min\{(-x)_+,\KM\}\le(-x)_+$ attains zero; any
$d(x)>(-x)_+$ gives a strictly positive value. Hence $d^{c}$ is the state-by-state
fossil minimizer, with \textsc{fcl} fossil $\equiv0$. Among the zero-fossil schedules
($d\le(-x)_+$), $d^{c}=\min\{(-x)_+,\KM\}$ is the one that runs the load to the full
available surplus, so it maximizes \textsc{fcl} revenue $R\,\E[d]$ subject to zero
fossil; it is therefore the zero-footprint contract the firm will actually operate
(any $d<d^{c}$ earns less surplus revenue for the same zero footprint, and not
operating at all, $d\equiv0$, forgoes the revenue that makes adoption worthwhile).
\end{proof}

So the surplus contract is the one a regulator who wants to keep the load's
benefits while erasing its fossil footprint would write: monetize the surplus, never
burn fossil to mine. Under it, \eqref{r1w3-eq:totfossil} collapses to consumer fossil,
\begin{equation}\label{r1w3-eq:fosc}
  \E[\fossil^{c}]=\int_0^{\KFc}\big(1-F(x;\KRc)\big)\,dx ,
\end{equation}
exactly the no-\textsc{fcl} form---but evaluated at the capacities the contract
induces, which is where the gain comes from.

%=====================================================================
\section{The decarbonization theorem (no shape condition)}\label{r1w3-sec:decarb}
%=====================================================================

Under the surplus contract the firm's program is
\begin{equation}\label{r1w3-eq:Pic}
  \Pi^{c}=p\,\E[\tilde D]
  +R\!\!\int_{-\KM}^{0}\!\!\!F\,dx
  -c\!\!\int_{0}^{\KF}\!\!(1-F)\,dx
  -(\bp)\!\!\int_{\KF}^{\Dbar}\!\!(1-F)\,dx
  -\CF\KF-\CR\KR-\CM\KM ,
\end{equation}
the \textsc{fcl} contributing the surplus revenue $R\,\E[\min\{(-\tilde x)_+,\KM\}]
=R\int_{-\KM}^{0}F\,dx$ and \emph{no} fossil cost. Its first-order conditions are
\begin{align}
  \partial_{\KM}\Pi^{c}&=R\,F(-\KM;\KR)-\CM=0
     &&\Longrightarrow\quad F(-\KMc;\KRc)=\tfrac{\CM}{R},\label{r1w3-eq:focM}\\
  \partial_{\KF}\Pi^{c}&=(\bp-c)\big(1-F(\KF;\KR)\big)-\CF=0
     &&\Longrightarrow\quad F(\KFc;\KRc)=\varphi^{\circ},\label{r1w3-eq:focF}\\
  \partial_{\KR}\Pi^{c}&=\underbrace{R\!\!\int_{-\KM}^{0}\!\!\FKR\,dx}_{>0\ \text{(adopted)},\ \text{surplus channel}}
     +\,c\!\!\int_{0}^{\KF}\!\!\FKR\,dx+(\bp)\!\!\int_{\KF}^{\Dbar}\!\!\FKR\,dx-\CR=0.\label{r1w3-eq:focR}
\end{align}
Two facts jump out. The fossil cutoff \eqref{r1w3-eq:focF} is \emph{identical} to the
no-\textsc{fcl} one \eqref{r1w3-eq:phio}: $F(\KFc;\KRc)=\varphi^{\circ}=F(\KFo;\KRo)$. And
the renewable condition \eqref{r1w3-eq:focR} is the no-\textsc{fcl} renewable condition plus
the strictly positive \emph{surplus channel} term $R\int_{-\KM}^{0}\FKR\,dx$. The
\textsc{fcl} is adopted ($\KMc>0$) when \eqref{r1w3-eq:focM} has an interior root; a
sufficient condition (using $F(0;\KRc)\ge F(0;\KRo)$, since $\KRc\ge\KRo$ and $F(0;\cdot)$
is increasing) is
\begin{equation}\label{r1w3-eq:adopt}
  \CM<R\,F(0;\KRo)\qquad\text{(adoption: the load is cheap enough to run on the surplus).}
\end{equation}

\begin{condition}[Single-peaked renewable problem]\label{r1w3-cond:reg}
The renewable marginal value (the left side of \eqref{r1w3-eq:focR} plus $\CR$) is strictly
decreasing in $\KR$ along the locus \eqref{r1w3-eq:focF}. This is the Week-2 reduced
second-order condition (the Schur/diagonal-dominance condition), here for the
surplus program; it holds on the regular region and at the power-law boundary it is
the knife-edge.
\end{condition}

The decarbonization theorem rests on a single monotone fact.

\begin{lemma}[More renewables, less fossil]\label{r1w3-lem:monofossil}
Along the fossil-cutoff locus $\KF(\KR)=F^{-1}(\varphi^{\circ};\KR)$, the utility's own
fossil output
$e(\KR)\eqdef\int_0^{(\KF(\KR))_+}\!\big(1-F(x;\KR)\big)\,dx$
is non-increasing in $\KR$, and strictly decreasing wherever $\KF(\KR)>0$:
\[
  e'(\KR)=\underbrace{(1-\varphi^{\circ})\,\KF'(\KR)}_{\le0}
  -\underbrace{\int_0^{\KF(\KR)}\!\FKR(x;\KR)\,dx}_{>0}\;<\;0,
  \qquad \KF'(\KR)=-\frac{\FKR(\KF;\KR)}{f(\KF;\KR)}\le0 .
\]
\end{lemma}
\begin{proof}
Differentiate $e$ by Leibniz, using $1-F(\KF;\KR)=1-\varphi^{\circ}$ at the moving upper
limit; $\KF'$ is the implicit derivative of $F(\KF(\KR);\KR)=\varphi^{\circ}$. Both terms
are $\le0$, the integral strictly so by $\FKR>0$: building renewables both lowers the
dispatchable-fossil cutoff and displaces fossil inside the served band. Where
$\KF(\KR)\le0$ the integral is empty and $e\equiv0$, so $e$ is non-increasing throughout.
\end{proof}

\begin{theorem}[Decarbonization, with no sign-determining shape condition]\label{r1w3-thm:decarb}
Under Condition~\ref{r1w3-cond:reg}, the surplus contract satisfies
\[
  \boxed{\;\KRc\ge\KRo\quad\text{and}\quad \E[\fossil^{c}]\le\E[\fossil^{\circ}],\;}
\]
both with \emph{strict} inequality exactly when the load is adopted \eqref{r1w3-eq:adopt}. The
carbon inequality is an immediate consequence of the renewable one: \emph{more renewables
mean less fossil}. The conclusion uses only the standing positivity $\FKR>0$ on the
interior support and the regularity Condition~\ref{r1w3-cond:reg} (which guarantees a
well-defined optimum and can fail, e.g.\ at the power-law knife-edge); given an optimum
and adoption, the \emph{sign} of the carbon change carries \emph{no} dependence on the
shape, tails, or skew of $F$. This is the sharp contrast with Week~2, where
(Remark~\ref{r1w3-rem:resolve}) the conditional shape of $F$ decided the sign; the surplus
contract removes that shape-dependence. (Magnitude still depends on shape; only the sign
is freed.)
\end{theorem}

\begin{proof}
\emph{Renewables ($\KRc\ge\KRo$).} If the load is not adopted ($\KMc=0$) the surplus
program is the no-\textsc{fcl} program and $\KRc=\KRo$. If it is adopted, evaluate the
renewable marginal value \eqref{r1w3-eq:focR} at $(\KRo,\KFo)$: by \eqref{r1w3-eq:focF} the fossil
cutoff there is $\KFo$, so its last two integrals are the no-\textsc{fcl} renewable
benefit, which equals $\CR$ at $\KRo$. The remaining surplus-channel term is strictly
positive,
\[
  \partial_{\KR}\Pi^{c}\big|_{(\KRo,\KFo)}=R\!\int_{-\KM}^{0}\!\FKR(x;\KRo)\,dx>0
  \qquad(\FKR>0\text{ on }(-\KM,0)),
\]
so the renewable marginal value exceeds $\CR$ at $\KRo$; by Condition~\ref{r1w3-cond:reg} it
is strictly decreasing, so the surplus optimum---its zero if interior, otherwise the
capacity bound $\bar K_R$---satisfies $\KRc>\KRo$.

\emph{Carbon ($\E[\fossil^{c}]\le\E[\fossil^{\circ}]$).} Under the surplus contract the
load burns no fossil (Prop.~\ref{r1w3-prop:surplusopt}), so total fossil is exactly the
utility's own output $e(\KRc)$, while the benchmark is $e(\KRo)$. By
Lemma~\ref{r1w3-lem:monofossil} $e$ is non-increasing and strictly decreasing while $\KF>0$;
with $\KRc\ge\KRo$ this gives $e(\KRc)\le e(\KRo)$, strict whenever $\KRc>\KRo$, i.e.\
whenever the load is adopted. That is the whole argument: the contract zeroes the load's
fossil and weakly raises renewables, and more renewables mean less fossil.\footnote{Verified
in \texttt{verify\_week3.py}: across symmetric, left- and right-skewed, bimodal, and
high-capacity-factor shapes, the surplus-channel term is strictly positive at $\KRo$,
$\KRc>\KRo$, and $\E[\fossil^{c}]<\E[\fossil^{\circ}]$ at every adopted optimum---the sign
never flips.}
\end{proof}

\begin{corollary}[Distributionally-robust decarbonization]\label{r1w3-cor:dro}
Let $\mathcal U$ be the ambiguity set of all residual-demand distributions satisfying the
standing positivity $\FKR>0$, Condition~\ref{r1w3-cond:reg}, and adoption \eqref{r1w3-eq:adopt}.
Then the surplus contract's carbon saving is worst-case \emph{nonnegative} over $\mathcal U$,
and strictly positive for each adopted law:
\[
  \inf_{F\in\mathcal U}\big(\E[\fossil^{\circ}]-\E[\fossil^{c}]\big)\;\ge\;0,
  \qquad \E[\fossil^{\circ}]-\E[\fossil^{c}]>0\ \text{ for every } F\in\mathcal U.
\]
So the regulator needs \emph{no knowledge of $F$}: whatever the true residual-demand law,
the surplus contract decarbonizes. Exogeneity of $F$ is not a limitation here---it is
precisely what the guarantee is robust to.
\end{corollary}
\begin{proof}
Theorem~\ref{r1w3-thm:decarb} gives $\E[\fossil^{\circ}]-\E[\fossil^{c}]>0$ for \emph{every}
$F\in\mathcal U$ (the proof uses only $\FKR>0$, Condition~\ref{r1w3-cond:reg}, and adoption).
A strict inequality holding for every member of a set gives $\inf\ge0$ (pointwise strictness
need not transfer to the infimum: as $F$ approaches marginal adoption, $K_M^{c}\to0$ and the
saving $\to0$, so the infimum can equal $0$). A uniform positive lower bound would need an
adoption margin bounded away from the threshold; we claim only the worst-case nonnegativity
and pointwise strict positivity, which is what the regulator's no-knowledge guarantee
requires.
\end{proof}

\begin{remark}[Why this resolves Week 2]\label{r1w3-rem:resolve}
Week~2 proved that under the naive contract the carbon sign is \emph{irreducibly
distributional}: at fixed equilibrium capacities the admissible conditional laws make
$\Delta$ range over an interval straddling zero, so no distribution-free statement is
possible and a fossil floor $(\KM-\KR)_+$ blocks reduction at interior optima.
Theorem~\ref{r1w3-thm:decarb} shows the obstruction was an artifact of the \emph{contract},
not the technology. The surplus contract sets the \textsc{fcl} fossil---and with it the
floor---to zero, and the only remaining channel is the renewable build-out, which has a
definite sign. The carbon question goes from ``it depends on the tails of $F$'' to
``it always falls when the load is adopted.''
\end{remark}

%=====================================================================
\section{The Pareto improvement and the exact saving}\label{r1w3-sec:triplewin}
%=====================================================================

\begin{corollary}[The surplus contract Pareto-dominates no load]\label{r1w3-cor:triplewin}
Under the hypotheses of Theorem~\ref{r1w3-thm:decarb}, relative to the no-\textsc{fcl}
benchmark the surplus contract \textup{(i)} strictly raises the utility's profit,
\textup{(ii)} holds the blackout probability \emph{exactly} at the no-load level, and
\textup{(iii)} strictly lowers carbon. It is thus a strict Pareto
improvement: more profit and less carbon at no reliability cost. The contract is
regulator-imposed, not firm-optimal: left to itself the firm would also run the load on
fossil headroom (the naive rule) and earn still more. But it strictly prefers operating
under the surplus contract to having no load, so participation is individually rational
and no subsidy is needed---the regulator gives up only the firm's extra
\emph{fossil-powered} computing profit, which is exactly the carbon-causing margin.
\end{corollary}

\begin{proof}
(iii) is Theorem~\ref{r1w3-thm:decarb}. (i): the surplus contract adds the option of
earning $R$ on otherwise-curtailed surplus at hardware cost $\CM$; by \eqref{r1w3-eq:adopt}
the first \textsc{fcl} unit's marginal surplus revenue $R\,F(0;\KRo)$ exceeds $\CM$, so
$\Pi^{c}>\Pi^{\circ}$ strictly (the firm can always replicate the benchmark by setting
$\KM=0$, and adoption is profitable). (ii): blackout occurs on $\{\tilde x>\KF\}$ with
probability $1-F(\KF;\KR)$. By the fossil cutoff \eqref{r1w3-eq:focF}, \emph{both} programs
target the same reliability, $F(\KFc;\KRc)=\varphi^{\circ}=F(\KFo;\KRo)$, so the
blackout probability is $1-\varphi^{\circ}$ in both---identical. The surplus load never
uses fossil, so it leaves the fossil-capacity trade-off, and hence the blackout target,
untouched; the larger $\KRc$ is spent on \emph{lowering fossil capacity}
($\KFc\le\KFo$, since the fossil cutoff falls in $\KR$ by Lemma~\ref{r1w3-lem:monofossil})
rather than on extra reliability. Reliability is therefore held,
not improved---a clean tie, which is what makes the result a Pareto
improvement.\footnote{Verified in \texttt{verify\_week3.py}: $\Pi^{c}>\Pi^{\circ}$ and
blackout$^{c}=$blackout$^{\circ}$ (to numerical tolerance) at every adopted optimum.}
\end{proof}

\begin{proposition}[Exact carbon saving and its monotonicity]\label{r1w3-prop:saving}
The carbon saving under the surplus contract is
\begin{equation}\label{r1w3-eq:exactsave}
  \Delta^{c}\eqdef\E[\fossil^{\circ}]-\E[\fossil^{c}]
  =\int_0^{\KFc}\!\!\big[F(x;\KRc)-F(x;\KRo)\big]\,dx
   +\int_{\KFc}^{\KFo}\!\!\big(1-F(x;\KRo)\big)\,dx \;>\;0 .
\end{equation}
It is increasing in the surplus channel: a larger adopted $\KMc$ (cheaper load, more
surplus) raises $\KRc$ and hence $\Delta^{c}$. In the cheap-renewable limit
$\KRc\to\infty$ the contract drives the utility to full decarbonization,
$\E[\fossil^{c}]\to0$, so $\Delta^{c}\to\E[\fossil^{\circ}]$.
\end{proposition}

\begin{proof}
In the interior case $\KFc>0$ the identity is the expansion of
$e(\KRo)-e(\KRc)$ used in Theorem~\ref{r1w3-thm:decarb}. Monotonicity is
Lemma~\ref{r1w3-lem:monofossil}: $\E[\fossil^{c}]=e(\KRc)$ is strictly decreasing in $\KRc$,
and a larger adopted $\KMc$ enlarges the surplus channel and raises $\KRc$, so
$\Delta^{c}=\E[\fossil^{\circ}]-e(\KRc)$ strictly increases. As $\KRc\to\infty$,
$F(x;\KRc)\to1$ for each fixed $x>0$ (renewables swamp demand), so by dominated
convergence $e(\KRc)\to0$ and $\Delta^{c}\to\E[\fossil^{\circ}]$ (full decarbonization).
\end{proof}

\begin{remark}[The contract incentivizes renewable build-out]
Proposition~\ref{r1w3-prop:saving} says the surplus contract does not merely fail to add
fossil---it actively pulls the optimal renewable capacity up, and in the cheap-renewable
regime all the way to full decarbonization. The flexible load, correctly contracted,
behaves like a subsidy for renewables financed by selling surplus power to computing.
This is the positive mechanism the paper has been after: not ``computing happens to be
green,'' but ``the right contract makes computing pay for the green build-out.''
\end{remark}

%=====================================================================
\section{The optimal-contract family and the carbon--profit frontier}\label{r1w3-sec:family}
%=====================================================================

So far we have two contracts: the Week-2 \emph{naive} rule (run the load whenever it
pays, on fossil too) and the Week-3 \emph{surplus} rule (run it on surplus only). They
are not two unrelated objects---they are the two ends of a single one-parameter family,
and seeing the whole family answers the question Week~3 left open: \emph{is surplus-only
the carbon-optimal contract, or could letting the load burn a little fossil decarbonize
even more, by buying a larger renewable build-out?}

\paragraph{The budget family.} For $\tau\in[0,\infty]$ let the regulator grant the load
a per-state \emph{fossil budget} $\tau$:
\begin{equation}\label{r1w3-eq:dtau}
  d_{\tau}(x)\eqdef\min\big\{(-x)_+ + \tau,\ \KM,\ (\KF-x)_+\big\}.
\end{equation}
At $\tau=0$ this is the surplus contract $d^{c}$ (\S\ref{r1w3-sec:optcontract}); as
$\tau\to\infty$ it is the naive contract $d^{*}$ (Week~2). The load's fossil is
$\Phi(\tau)\eqdef\E[(d_{\tau}(\tilde x)-(-\tilde x)_+)_+]$, with $\Phi(0)=0$ and $\Phi$
increasing. For each $\tau$ the firm re-optimizes all three capacities
$(\KR(\tau),\KF(\tau),\KM(\tau))$; write $\Pi(\tau)$ for its profit and
\begin{equation}\label{r1w3-eq:Ctau}
  C(\tau)=\underbrace{G\big(\KR(\tau),\KF(\tau)\big)}_{\text{consumer fossil}}+\Phi(\tau),
  \qquad G(\KR,\KF)\eqdef\int_0^{\KF}\!\big(1-F(x;\KR)\big)\,dx ,
\end{equation}
for total carbon: consumer fossil $G$ at the induced capacities plus the load's direct
footprint $\Phi(\tau)$. At $\tau=0$, $\Phi=0$ and the fossil cutoff is $\varphi^{\circ}$,
so $C(0)=e(\KRc)$ recovers the surplus value of Lemma~\ref{r1w3-lem:monofossil}; for $\tau>0$
the load draws fossil, so the firm builds \emph{more} fossil capacity ($\KF(\tau)$ rises)
and $G$ must be evaluated at that larger $\KF$, not on the $\varphi^{\circ}$ locus.

\begin{proposition}[Profit rises with the budget]\label{r1w3-prop:profitmono}
$\Pi(\tau)$ is non-decreasing in $\tau$, and strictly increasing wherever the budget
binds.
\end{proposition}
\begin{proof}
For fixed capacities, enlarging $\tau$ enlarges the dispatch upper bound in
\eqref{r1w3-eq:dtau} pointwise, so the firm can dispatch weakly more; each extra unit earns
$R$ on surplus or $R-c>0$ on fossil (Assumption~\ref{r1w3-ass:reg}), both positive, so
fixed-capacity profit is non-decreasing in $\tau$. Taking the max over capacities
preserves this. Where the budget binds the firm strictly dispatches more, so the
increase is strict; where instead the cap $\KM$ or the headroom $(\KF-x)_+$ binds, $\Pi$
is flat in $\tau$ (a larger budget changes nothing).
\end{proof}

Write $K(\tau)=(\KR(\tau),\KF(\tau),\KM(\tau))$ for the firm's $\tau$-optimal capacities
and define the \emph{channel leverage}
\begin{equation}\label{r1w3-eq:leverage}
  \Lambda(\tau)\eqdef\frac{G(K(0))-G(K(\tau))}{\Phi(\tau)}\qquad(\tau>0),
\end{equation}
the consumer fossil that the budget's induced capacity change saves, per unit of the
direct footprint $\Phi(\tau)$ it burns.

\begin{theorem}[The carbon--profit frontier; surplus-only minimizes carbon]\label{r1w3-thm:frontier}
Over the family \eqref{r1w3-eq:dtau}:
\begin{enumerate}[leftmargin=2em,itemsep=1pt]
\item[\textup{(i)}] $\Pi(\tau)$ is non-decreasing (Proposition~\ref{r1w3-prop:profitmono}), so
$(\Pi(\tau),C(\tau))$ traces a monotone frontier: a larger fossil budget weakly buys firm
profit;
\item[\textup{(ii)}] \textbf{(exact reduction)} $C(\tau)-C(0)=\Phi(\tau)\,\big(1-\Lambda(\tau)\big)$;
hence surplus-only ($\tau=0$) is the carbon-optimal member of the family \emph{if and only
if} $\Lambda(\tau)\le1$ for every $\tau>0$---no budget's renewable build-out saves more
fossil than the budget directly burns;
\item[\textup{(iii)}] the carbon-minimizing contract subject to a firm-profit floor
$\bar\Pi\in[\Pi^{c},\Pi^{*}]$ is $\tau^{*}=\inf\{\tau:\Pi(\tau)\ge\bar\Pi\}$; it is the
unique carbon-minimizer whenever $\Lambda\le1$ holds strictly, and $\tau^{*}=0$
(surplus-only) when the floor is slack ($\bar\Pi\le\Pi^{c}$).
\end{enumerate}
The condition $\Lambda(\tau)\le1$ holds in every shape and regime we scan---indeed
$\Lambda(\tau)<0$, since the budget-induced capacities \emph{raise} consumer fossil
(a larger $\KF$ to serve the load), so there is no offset at all and $C(\tau)>C(0)$
strictly.
\end{theorem}
\begin{proof}
(ii) is an exact identity: by \eqref{r1w3-eq:Ctau}, $C(\tau)=G(K(\tau))+\Phi(\tau)$ and
$C(0)=G(K(0))$ (as $\Phi(0)=0$), so
$C(\tau)-C(0)=\Phi(\tau)-\big[G(K(0))-G(K(\tau))\big]=\Phi(\tau)\big(1-\Lambda(\tau)\big)$.
Since $\Phi(\tau)>0$ for $\tau>0$, $C(\tau)\ge C(0)\Leftrightarrow\Lambda(\tau)\le1$---no
approximation, no unproven capacity-response sign. (i) is Proposition~\ref{r1w3-prop:profitmono}.
(iii) The floor-feasible set is $\{\tau:\Pi(\tau)\ge\bar\Pi\}=[\tau^{*},\infty)$ ($\Pi$
non-decreasing); under $\Lambda\le1$, $C$ is non-decreasing so the left end $\tau^{*}$
minimizes carbon, uniquely where the inequality is strict; $\bar\Pi\le\Pi^{c}=\Pi(0)$
gives $\tau^{*}=0$.
\end{proof}

\begin{remark}[Why $\Lambda<1$, and what would overturn it]
The leverage $\Lambda$ is the sharp knob. Differentiating \eqref{r1w3-eq:Ctau},
$dC/d\tau=\Phi'(\tau)+G_{\KF}\KF'(\tau)+G_{\KR}\KR'(\tau)$ with
$G_{\KF}=1-F(\KF)\ge0$, $G_{\KR}=-\int_0^{\KF}\FKR<0$: two forces push carbon up
(the footprint $\Phi'>0$ and the extra consumer fossil from the larger $\KF$ the load
needs), and the \emph{only} way to get $\Lambda>1$ is a renewable channel $G_{\KR}\KR'$
strong enough---$\KR'(\tau)$ both positive and large---to overturn both. That is exactly
what the theorem isolates as the single thing to check, and across every shape and regime
$\Lambda<0<1$: fossil-running does not even raise $\KR$ on net here, it raises $\KF$, so
the ``financing'' story fails and carbon rises unambiguously. The result is thus a
\emph{proved} equivalence resting on one interpretable, verified scalar---not a
distribution-free assertion about $\KR'(\tau)$, which we do not claim.
\end{remark}

\begin{remark}[No free lunch, and the unification]
Within this budget family, the carbon question reduces to a single monotone trade-off:
every megawatt of fossil the load is permitted to burn raises emissions, because---in
every case we compute---the renewable build-out that fossil-running finances does not
pay for the fossil it costs. So the carbon-optimal contract in the family is the extreme
one, zero fossil for the load, and Weeks~2 and~3 are not rival models but the two ends of
one frontier: the naive rule is the profit-maximizing, carbon-worst corner, the surplus
rule the carbon-best corner. We characterize this positive menu---the carbon and firm
profit of each contract; \emph{which} point a planner should select (the welfare and
private-vs-social question, and the blackout--emissions trade-off) is a
reliability-and-welfare matter left to the companion workstream, not an A1 claim. (We
also claim optimality \emph{within the family}, not over all conceivable contracts; the
broader bilevel question is flagged in \S\ref{r1w3-sec:optcontract}.) This rounds out the
optimal-contract account of the flexible load's carbon effect.
\end{remark}

%=====================================================================
\section{Numerical validation}\label{r1w3-sec:num}

All claims are checked in \texttt{verify\_week3.py} against adaptive quadrature of the
profit integrals, on non-power-law Beta scale mixtures
$\tilde x=\tilde D-\KR\tilde S$ spanning symmetric, left-skewed, right-skewed, and
bimodal demand. The suite confirms, at every adopted optimum and across the family:
the surplus contract burns zero \textsc{fcl} fossil (Prop.~\ref{r1w3-prop:surplusopt}); the
surplus-channel term in \eqref{r1w3-eq:focR} is strictly positive at $\KRo$, giving
$\KRc>\KRo$; the fossil output $e(\KR)$ is strictly decreasing
(Lemma~\ref{r1w3-lem:monofossil}); and the strict carbon fall
$\E[\fossil^{c}]<\E[\fossil^{\circ}]$ holds with the \emph{same sign for every shape in
the family} (Theorem~\ref{r1w3-thm:decarb}); the Pareto improvement $\Pi^{c}>\Pi^{\circ}$
with blackout held equal, blackout$^{c}=$blackout$^{\circ}=1-\varphi^{\circ}$
(Cor.~\ref{r1w3-cor:triplewin}); and the exact saving \eqref{r1w3-eq:exactsave} matches the direct
quadrature. As a contrast the suite also
reports the Week-2 naive-rule fossil at the same primitives, exhibiting the floor that
the contract removes. The suite reports a single pass count.
